\documentclass[a4paper]{article}
\usepackage[english]{babel}
\usepackage{amsmath}
\usepackage{amsfonts}
\usepackage{amssymb}
\usepackage{bbold}
\usepackage{graphicx}
\usepackage{hyperref}
\usepackage{amsthm}
\newtheorem{theorem}{Theorem}
\newtheorem{lemma}{Lemma}
\usepackage{algorithm}
\usepackage{algpseudocode}
\usepackage{booktabs}

\setcounter{MaxMatrixCols}{20}

\begin{document}
\pagenumbering{arabic}

\Large
\begin{center}
\textbf{Finite Difference Method for Hyperbolic Equations\\}

\hspace{10pt}

\large
Julio A. Medina\\
\hspace{10pt}
\small
Universidad de San Carlos\\
School of Physical Sciences and Mathematics\\
Master's Program in Physics\\
\href{mailto:julioantonio.medina@gmail.com}{julioantonio.medina@gmail.com}\\
\end{center}

\hspace{10pt}

\normalsize
\section{Hyperbolic Partial Differential Equations}

The hyperbolic partial differential equation, also known as the wave equation, is given by
\begin{equation}\label{eq::hyperbolic_partial_diff}
\frac{\partial^2 u}{\partial t^2}(x,t)-\alpha^2\frac{\partial^2 u}{\partial x^2}(x,t)=0,\qquad 0<x<l,\qquad t>0.
\end{equation}
It is subject to the boundary conditions
\begin{equation}
u(0,t)=u(l,t)=0,\qquad \text{for }t>0,
\end{equation}
and the initial conditions
\begin{equation}
u(x,0)=f(x),\qquad \frac{\partial u}{\partial t}(x,0)=g(x),\qquad \text{for }0\leq x\leq l,
\end{equation}
where $\alpha$ is a constant determined by the physical conditions of the problem. The wave equation is ubiquitous throughout physics.

\subsection{Finite Difference Method for the Wave Equation}

As in the finite-difference methods used to solve parabolic and elliptic equations, an integer $m>0$ is selected to define lattice points along the $x$-axis, with $h=l/m$. A temporal step size $k>0$ is also selected. The grid points $(x_i,t_j)$ are defined by
\begin{equation}
x_i=ih,\qquad t_j=jk,
\end{equation}
for $i=0,1,2,\ldots,m$ and $j=0,1,2,\ldots$.

At every interior grid point, namely, every point that does not lie on the grid boundary, the wave equation in Eq.~\ref{eq::hyperbolic_partial_diff} becomes
\begin{equation}\label{eq::wave_equation}
\frac{\partial^2 u}{\partial t^2}(x_i,t_j)-\alpha^2\frac{\partial^2 u}{\partial x^2}(x_i,t_j)=0.
\end{equation}
The finite-difference method is obtained by approximating the second partial derivatives with centered-difference quotients; see Ref.~\cite{Burden}. These approximations are
\begin{equation}\label{eq::difference1}
\frac{\partial^2 u}{\partial t^2}(x_i,t_j)=\frac{u(x_i,t_{j+1})-2u(x_i,t_j)+u(x_i,t_{j-1})}{k^2}-\frac{k^2}{12}\frac{\partial^4 u}{\partial t^4}(x_i,\mu_j),
\end{equation}
where $\mu_j\in(t_{j-1},t_{j+1})$, and
\begin{equation}\label{eq::difference2}
\frac{\partial^2 u}{\partial x^2}(x_i,t_j)=\frac{u(x_{i+1},t_j)-2u(x_i,t_j)+u(x_{i-1},t_j)}{h^2}-\frac{h^2}{12}\frac{\partial^4 u}{\partial x^4}(\xi_i,t_j),
\end{equation}
where $\xi_i\in(x_{i-1},x_{i+1})$. Substituting Eqs.~\ref{eq::difference1} and \ref{eq::difference2} into Eq.~\ref{eq::wave_equation} gives
\begin{equation}
\begin{aligned}
&\frac{u(x_i,t_{j+1})-2u(x_i,t_j)+u(x_i,t_{j-1})}{k^2}-\alpha^2\frac{u(x_{i+1},t_j)-2u(x_i,t_j)+u(x_{i-1},t_j)}{h^2}\\
&\qquad=\frac{1}{12}\left[k^2\frac{\partial^4u}{\partial t^4}(x_i,\mu_j)-\alpha^2h^2\frac{\partial^4u}{\partial x^4}(\xi_i,t_j)\right].
\end{aligned}
\end{equation}
Neglecting the error term involving the fourth derivatives in time and space yields the difference equation
\begin{equation}\label{eq::difference_equation}
\frac{w_{i,j+1}-2w_{i,j}+w_{i,j-1}}{k^2}-\alpha^2\frac{w_{i+1,j}-2w_{i,j}+w_{i-1,j}}{h^2}=0.
\end{equation}
Defining
\begin{equation}
\lambda=\frac{\alpha k}{h},
\end{equation}
Eq.~\ref{eq::difference_equation} can be rewritten as
\begin{equation}\label{eq::difference_equation_2}
w_{i,j+1}-2w_{i,j}+w_{i,j-1}-\lambda^2w_{i+1,j}+2\lambda^2w_{i,j}-\lambda^2w_{i-1,j}=0.
\end{equation}
Solving for the approximation at the next time level gives
\begin{equation}\label{eq::difference_equation_3}
w_{i,j+1}=2(1-\lambda^2)w_{i,j}+\lambda^2\left(w_{i+1,j}+w_{i-1,j}\right)-w_{i,j-1}.
\end{equation}
This equation is valid for $i=1,2,\ldots,m-1$ and $j=1,2,\ldots$. The boundary conditions give
\begin{equation}\label{eq::condition_1}
w_{0,j}=w_{m,j}=0,\qquad \forall j=1,2,3,\ldots,
\end{equation}
and the initial displacement implies
\begin{equation}
w_{i,0}=f(x_i),\qquad \forall i=1,2,\ldots,m-1.
\end{equation}
For the standard explicit scheme, stability requires the Courant--Friedrichs--Lewy condition $\lambda\leq 1$.

The recurrence can be written conveniently in matrix form:
\begin{equation}\label{eq::matrix_difference_equation}
\begin{bmatrix}
w_{1,j+1}\\
w_{2,j+1}\\
w_{3,j+1}\\
\vdots\\
w_{m-1,j+1}
\end{bmatrix}
=
\begin{bmatrix}
2(1-\lambda^2)&\lambda^2&0&\dots&0\\
\lambda^2&\ddots&\ddots&\ddots&\vdots\\
0&\ddots&\ddots&\ddots&0\\
\vdots&\ddots&\ddots&\ddots&\lambda^2\\
0&\dots&0&\lambda^2&2(1-\lambda^2)
\end{bmatrix}
\begin{bmatrix}
w_{1,j}\\
w_{2,j}\\
w_{3,j}\\
\vdots\\
w_{m-1,j}
\end{bmatrix}
-
\begin{bmatrix}
w_{1,j-1}\\
w_{2,j-1}\\
w_{3,j-1}\\
\vdots\\
w_{m-1,j-1}
\end{bmatrix}.
\end{equation}

Equations~\ref{eq::difference_equation} and \ref{eq::difference_equation_3} show that the $(j+1)$th time step requires knowledge of the solutions at time levels $j$ and $j-1$. This creates an initialization problem: the values at $j=0$ are given by the initial displacement, but the values at $j=1$, which are needed to compute $w_{i,2}$, must be obtained from the initial velocity condition
\begin{equation}
\frac{\partial u}{\partial t}(x,0)=g(x),\qquad 0\leq x\leq l.
\end{equation}
One approach is to replace $\partial u/\partial t$ with a forward-difference approximation:
\begin{equation}
\frac{\partial u}{\partial t}(x_i,0)=\frac{u(x_i,t_1)-u(x_i,0)}{k}-\frac{k}{2}\frac{\partial^2u}{\partial t^2}(x_i,\widetilde{\mu}_i),
\end{equation}
for some $\widetilde{\mu}_i\in(0,t_1)$. Solving for $u(x_i,t_1)$ gives
\begin{equation}
\begin{aligned}
u(x_i,t_1)&=u(x_i,0)+k\frac{\partial u}{\partial t}(x_i,0)+\frac{k^2}{2}\frac{\partial^2u}{\partial t^2}(x_i,\widetilde{\mu}_i)\\
&=u(x_i,0)+kg(x_i)+\frac{k^2}{2}\frac{\partial^2u}{\partial t^2}(x_i,\widetilde{\mu}_i).
\end{aligned}
\end{equation}
Omitting the truncation error gives the approximation
\begin{equation}
w_{i,1}=w_{i,0}+kg(x_i),\qquad \forall i=1,2,\ldots,m-1.
\end{equation}
However, this approximation has truncation error $O(k)$, whereas the centered recurrence in Eq.~\ref{eq::difference_equation_3} is second-order accurate in time.

\subsection{Improving the Initial Approximation}

To obtain a more accurate approximation for $u(x_i,t_1)$, expand it in a second-order Maclaurin polynomial in $t$:
\begin{equation}
u(x_i,t_1)=u(x_i,0)+k\frac{\partial u}{\partial t}(x_i,0)+\frac{k^2}{2}\frac{\partial^2u}{\partial t^2}(x_i,0)+\frac{k^3}{6}\frac{\partial^3u}{\partial t^3}(x_i,\widehat{\mu}_i),
\end{equation}
for some $\widehat{\mu}_i\in(0,t_1)$. If $f''$ exists, then
\begin{equation}
\frac{\partial^2u}{\partial t^2}(x_i,0)=\alpha^2\frac{\partial^2u}{\partial x^2}(x_i,0)=\alpha^2\frac{\partial^2f}{\partial x^2}(x_i)=\alpha^2f''(x_i).
\end{equation}
Therefore,
\begin{equation}
u(x_i,t_1)=u(x_i,0)+kg(x_i)+\frac{\alpha^2k^2}{2}f''(x_i)+\frac{k^3}{6}\frac{\partial^3u}{\partial t^3}(x_i,\widehat{\mu}_i).
\end{equation}
This gives the approximation
\begin{equation}\label{eq::aprox_1}
w_{i,1}=w_{i,0}+kg(x_i)+\frac{\alpha^2k^2}{2}f''(x_i),
\end{equation}
with a local truncation error of order $O(k^3)$.

If $f\in C^4[0,l]$ but $f''(x_i)$ is not directly available, the centered-difference formula
\begin{equation}\label{eq::second_derivative_diff}
f''(x_0)=\frac{f(x_0-h)-2f(x_0)+f(x_0+h)}{h^2}-\frac{h^2}{12}f^{(4)}(\xi)
\end{equation}
can be used for some $\xi$ satisfying $x_0-h<\xi<x_0+h$. Thus,
\begin{equation}
f''(x_i)=\frac{f(x_{i+1})-2f(x_i)+f(x_{i-1})}{h^2}-\frac{h^2}{12}f^{(4)}(\widetilde{\xi}_i),
\end{equation}
for some $\widetilde{\xi}_i\in(x_{i-1},x_{i+1})$. It follows that
\begin{equation}
u(x_i,t_1)=u(x_i,0)+kg(x_i)+\frac{k^2\alpha^2}{2h^2}\left[f(x_{i+1})-2f(x_i)+f(x_{i-1})\right]+O(k^3+h^2k^2).
\end{equation}
Using $\lambda=\alpha k/h$ gives
\begin{equation}
\begin{aligned}
u(x_i,t_1)&=u(x_i,0)+kg(x_i)+\frac{\lambda^2}{2}\left[f(x_{i+1})-2f(x_i)+f(x_{i-1})\right]+O(k^3+h^2k^2)\\
&=(1-\lambda^2)f(x_i)+\frac{\lambda^2}{2}f(x_{i+1})+\frac{\lambda^2}{2}f(x_{i-1})+kg(x_i)+O(k^3+h^2k^2).
\end{aligned}
\end{equation}
The improved initialization formula is therefore
\begin{equation}\label{eq::aprox_improvement}
w_{i,1}=(1-\lambda^2)f(x_i)+\frac{\lambda^2}{2}f(x_{i+1})+\frac{\lambda^2}{2}f(x_{i-1})+kg(x_i).
\end{equation}
It can be used to compute $w_{i,1}$ for every $i=1,2,\ldots,m-1$. The subsequent approximations are obtained from the system in Eq.~\ref{eq::matrix_difference_equation}.

\section{Finite Difference Algorithm for the Wave Equation}

The development in the preceding section provides the tools needed to construct an algorithm for hyperbolic partial differential equations of the form in Eq.~\ref{eq::hyperbolic_partial_diff}. The procedure is as follows:
\begin{itemize}
\item Construct the matrix system in Eq.~\ref{eq::matrix_difference_equation}.
\item Use the improved initial approximation in Eq.~\ref{eq::aprox_improvement} to compute $w_{i,1}$.
\item Apply the matrix recurrence for the required number of time steps.
\end{itemize}

\subsection{Algorithm}

The following pseudocode implements the explicit centered finite-difference scheme together with the improved first time step.
\begin{algorithm}[H]
\caption{Explicit finite-difference method for the wave equation}\label{alg::finite_difference_hyperbolic}
\begin{algorithmic}[1]
\Function{finite\_difference\_method\_hyperbolic}{$l,T,\alpha,N,m,f,g$}
\State $h\gets l/m$
\State $k\gets T/N$
\State $\lambda\gets \alpha k/h$
\State $x\gets \operatorname{linspace}(0,l,m+1)$
\State $t\gets \operatorname{linspace}(0,T,N+1)$
\State $A\gets \operatorname{zeros}(m-1,m-1)$
\State $w_0\gets f(x[1:-1])$
\State $w_1\gets(1-\lambda^2)f(x[1:-1])+\frac{\lambda^2}{2}\left(f(x[2:])+f(x[0:-2])\right)+kg(x[1:-1])$
\For{$i\gets0$ to $m-2$}
\State $A[i,i]\gets2(1-\lambda^2)$
\If{$i<m-2$}
\State $A[i,i+1]\gets\lambda^2$
\State $A[i+1,i]\gets\lambda^2$
\EndIf
\EndFor
\For{each $t_j$ in $t[2:]$}
\State $w_j\gets Aw_1-w_0$
\State $w_0\gets w_1$
\State $w_1\gets w_j$
\EndFor
\State \Return $A,w_j,w_0,x$
\EndFunction
\end{algorithmic}
\end{algorithm}

The corresponding implementation is available in the repository at \url{https://github.com/Julio-Medina/Finite_Difference_Method_Hyperbolic}.

\subsection{Example}

Use the finite-difference algorithm in Algorithm~\ref{alg::finite_difference_hyperbolic}, with $h=0.1$ and $k=0.05$, to approximate the solution of the wave equation
\begin{equation}
\frac{\partial^2u}{\partial t^2}(x,t)-4\frac{\partial^2u}{\partial x^2}(x,t)=0,\qquad 0<x<1,
\qquad t>0,
\end{equation}
subject to the boundary conditions
\begin{equation*}
u(0,t)=u(1,t)=0,\qquad t>0,
\end{equation*}
and the initial conditions
\begin{equation}
u(x,0)=\sin(\pi x),\qquad \frac{\partial u}{\partial t}(x,0)=0,
\qquad 0\leq x\leq1.
\end{equation}
The analytical solution is
\begin{equation}
u(x,t)=\sin(\pi x)\cos(2\pi t).
\end{equation}

For $m=10$, $N=20$, and $\alpha=2$, the grid spacings are $h=0.1$ and $k=0.05$, and the Courant number is $\lambda=1$. The numerical and analytical values at $t=1$ are compared in Table~\ref{Tab::table_1}.

\begin{table}[ht]
\centering
\caption{Numerical and analytical solutions at $t=1$.}
\label{Tab::table_1}
\begin{tabular}{lrrrr}
\toprule
$i$&$x_i$&$w_{i,20}$&$u(x_i,1)$&$\lvert u(x_i,1)-w_{i,20}\rvert$\\
\midrule
0&0.0&0.000000&0.000000e+00&0.000000e+00\\
1&0.1&0.309017&3.090170e-01&5.551115e-17\\
2&0.2&0.587785&5.877853e-01&2.220446e-16\\
3&0.3&0.809017&8.090170e-01&0.000000e+00\\
4&0.4&0.951057&9.510565e-01&3.330669e-16\\
5&0.5&1.000000&1.000000e+00&0.000000e+00\\
6&0.6&0.951057&9.510565e-01&2.220446e-16\\
7&0.7&0.809017&8.090170e-01&0.000000e+00\\
8&0.8&0.587785&5.877853e-01&0.000000e+00\\
9&0.9&0.309017&3.090170e-01&5.551115e-17\\
10&1.0&0.000000&1.224647e-16&1.224647e-16\\
\bottomrule
\end{tabular}
\end{table}

\section{Conclusions}

The finite-difference method for hyperbolic partial differential equations with boundary conditions and initial displacement and velocity data has a relatively straightforward implementation. The method is based on constructing a tridiagonal matrix recurrence and advancing the numerical solution in time. The improved initialization formula preserves the second-order temporal accuracy of the centered scheme.

As shown in Table~\ref{Tab::table_1}, the numerical solution agrees with the analytical solution to approximately machine precision for the selected example. This exceptionally small error is a special property of the chosen parameters: $\lambda=1$, and the initial condition is a discrete sine eigenmode whose numerical frequency matches the exact frequency. In general, the centered scheme is second-order accurate in both time and space, provided the solution is sufficiently smooth and the CFL stability condition is satisfied.

\begin{thebibliography}{99}

\bibitem{Burden} Richard L. Burden and J. Douglas Faires, \textit{Numerical Analysis}, 9th ed., Brooks/Cole, Cengage Learning, ISBN 978-0-538-73351-9.

\bibitem{Medina} Julio Medina, \textit{Finite Difference Method for Elliptic Equations}, \url{https://github.com/Julio-Medina/Finite_Difference_Method}.

\bibitem{Varga} Richard S. Varga, \textit{Matrix Iterative Analysis}, 2nd ed., Springer, DOI: 10.1007/978-3-642-05156-2.

\end{thebibliography}
\end{document}
